\subsection{A determinant obstruction without enumerating sphericalizations}
\label{sec:signed-determinant}
The field and divisibility conditions can be combined before determining any
signed multiplication table. This gives a general-categorification obstruction without a bounded signed-table search.

\begin{proposition}[Signed determinant bound]\label{prop:signed-determinant}
Let $R$ be a commutative self-dual based ring of rank $r$, with fusion
matrices $N_0=I,N_1,\ldots,N_{r-1}$, and put
\[
 H_i=\left(\sum_jN_j^2\right)_{ii}.
\]
Suppose that the field generated by all formal codegrees of $R$ has
degree at least $r$. For each distinct irreducible factor $P$ of the formal-codegree
polynomial put $e_P=\deg P$ and $n_P=|P(0)|$, and define
\[
 B=\prod_p p^{\left\lceil r\max_P
          \{v_p(n_P)/e_P\}\right\rceil}.
\]
If $R$ admits a fusion categorification over $\mathbb C$, then
\[
 B\leq\prod_{i=0}^{r-1}H_i.
\]
\end{proposition}
\begin{proof}
Form the signed quotient $A$ of the sphericalization and let $\chi$ be its
spherical dimension character. Write $L=\mathbb Q(\chi(b_i):0\leq i<r)$.
The codegree-field condition gives $[L:\mathbb Q]\geq r$. The character
map $A_{\mathbb Q}\to L$ is surjective: its image is a finite-dimensional
integral domain over $\mathbb Q$, hence a field. Since
$\dim_{\mathbb Q}A_{\mathbb Q}=r$, this map is an isomorphism.
In particular $[L:\mathbb Q]=r$.

Let $T_i$ denote signed multiplication by $b_i$. These are real symmetric
matrices, and the signed Casimir
\[
 Q_T=\sum_iT_i^2=I+\sum_{i>0}T_i^2
\]
is positive definite. Its value under $\chi$ is
$\delta=\sum_i\chi(b_i)^2$, and consequently
$\det Q_T=N_{L/\mathbb Q}(\delta)$.
For every original formal codegree $f$, the ratio $\delta/f$ is an
algebraic integer by \cite[Proposition~\ref{B-prop:signed}]{Finite}. Taking norms in $L$ gives
\[
 v_p(\det Q_T)\geq r\,v_p(n_P)/e_P
 \qquad(f\text{ a root of }P).
\]
The determinant is a positive integer, so it is a multiple of $B$.
On the other hand $|(T_j)_{ik}|\leq(N_j)_{ik}$ implies
\[
 (Q_T)_{ii}=\sum_{j,k}(T_j)_{ik}^2\leq H_i.
\]
Hadamard's determinant inequality therefore gives
$B\leq\det Q_T\leq\prod_i(Q_T)_{ii}\leq\prod_iH_i$.
\end{proof}

\begin{corollary}\label{cor:two-quadratic-determinant}
Suppose a self-dual rank-four ring has two irreducible quadratic
formal-codegree factors with distinct quadratic fields and norms $a,b$.
If it is categorifiable over $\mathbb C$, then
\[
 \operatorname{lcm}(a,b)^2\leq H_0H_1H_2H_3.
\]
\end{corollary}
\begin{proof}
The two quadratic fields generate a biquadratic field. Their norms give
$B=\operatorname{lcm}(a,b)^2$ in Proposition~\ref{prop:signed-determinant}.
\end{proof}

\begin{remark}[Sufficiency for $d$-numbers]
The normalized norm inequalities are not generally sufficient for
algebraic integrality of a quotient. They are sufficient when both
numerator $u$ and denominator $v$ are $d$-numbers. Enlarge $\Q(u,v)$
to a normal field $L$. The $d$-number property makes their valuations
constant on all primes over a fixed rational prime $p$. If the common
ramification index is $e$, those valuations are
$e\,v_p(n_u)/[\Q(u):\Q]$ and
$e\,v_p(n_v)/[\Q(v):\Q]$, respectively. The normalized inequalities
therefore make $u/v$ integral at every finite prime. Only the necessary
direction is used in the finite-classification checker.
\end{remark}
